\documentclass[a4paper,11pt]{article}

\usepackage{revy}
\usepackage[utf8]{inputenc}
\usepackage[T1]{fontenc}
\usepackage[danish]{babel}
\usepackage{amssymb}


\revyname{MatematikRevyen}
\revyyear{2026}
% HUSK AT OPDATERE VERSIONSNUMMER
\version{0.1}
\eta{$n$ minutter}
\status{Færdig}

\title{Par-tal-lederdebat}
\author{Louie '24, Thais '22 og Philip '24}

\begin{document}
\maketitle

\begin{roles}
\role{Clement}[Skuespiller] Debatleder
\role{\ensuremath{\emptyset}}[Skuespiller] Liste \ensuremath{\emptyset}, den tomme mængde. Siger ingenting
\role{\ensuremath{\mathbb{N}}}[Skuespiller] Lalleglade. Siger til alle de andre, at deres holdninger er så ringe
\role{\ensuremath{\mathbb{N}_0}}[Skuespiller] Lille fringe-parti med hjemmelavet podie. Mener selv, han er \ensuremath{\mathbb{N}}, bare bedre
\role{\ensuremath{\mathbb{Z}}}[Skuespiller] Spytter ikke facts, skøre
\role{\ensuremath{\mathbb{Q}}}[Skuespiller] Det rationelle parti
\role{\ensuremath{\mathbb{R}}}[Skuespiller] Seriøse og lidt arrogante. Føler sig bedre end de andre: ``Vi er bare mere fuldstændige end jer''
\role{\ensuremath{\mathbb{C}}}[Skuespiller] Nyt parti, lidt over det hele
\role{Ninjaer}[Scenefolk] Fjerner \ensuremath{\mathbb{N}_0}
\role{Publikum} Publikum
\end{roles}

\begin{props}
\prop{Podier}[Person, der skaffer] Et til hvert parti, med partiets mængdesymbol på
\prop{Hjemmelavet podie}[Person, der skaffer] Til \ensuremath{\mathbb{N}_0}, skal tydeligt se hjemmelavet ud
\end{props}


\begin{sketch}

\scene{Sceneopstilling fra venstre: \ensuremath{\emptyset}, \ensuremath{\mathbb{N}}, (\ensuremath{\mathbb{N}_0}), \ensuremath{\mathbb{Z}}, \ensuremath{\mathbb{Q}}, \ensuremath{\mathbb{R}}, \ensuremath{\mathbb{C}}. Alle partierne står bag hver deres podie med deres mængdesymbol på. Clement står ved siden af og introducerer debatten til publikum.}

\says{Clement} Velkommen til aftenens par-tal-lederdebat.

\scene{Venter på klap fra publikum}

\says{Clement} Det her kommer jo til at afgøre, hvem der bliver den næste premierminister.

\says{Hvor ville I placere jer selv på aksen?}

\says{\ensuremath{\mathbb{N}}} Altså vi ligger jo skarpt til højre for midten.

\says{\ensuremath{\mathbb{N}}} Vi har den helt positive tilgang til verdens problemer.

\says{\ensuremath{\mathbb{N}}} Vi er slet ikke negative som alle jer andre.

\says{\ensuremath{\mathbb{Z}}} I kan ikke bare lalle rundt og ignorere alt det negative.

\says{\ensuremath{\mathbb{N}}} Jo, vi ka så.

\says{\ensuremath{\mathbb{Q}}} Men \ensuremath{\mathbb{Z}}, du skal heller ikke komme for godt i gang.

\says{\ensuremath{\mathbb{Q}}} Der er virkelig mange steder, hvor I slet ikke mener noget.

\says{\ensuremath{\mathbb{Z}}} Vi går ikke ned i de helt små detaljer, men vi har holdninger til det hele.

\says{\ensuremath{\mathbb{R}}} Men \ensuremath{\mathbb{Q}}, der er jo også ret mange steder, hvor I ikke mener noget.

\says{\ensuremath{\mathbb{Q}}} Det er så lille en forskel, vi ligger i hvert fald tæt nok på.

\says{Clement} Hvad med dig \ensuremath{\mathbb{C}}, du har ikke sagt noget endnu?

\says{\ensuremath{\mathbb{C}}} Vi har faktisk ikke lyst til at placere os på den akse, du nu har givet os.

\says{\ensuremath{\mathbb{C}}} Vi ser ikke på tingene så endimensionelt som de andre partier.

\says{\ensuremath{\mathbb{C}}} Matematikken er alt for kompleks til at reducere.

\says{Clement} Empty set får sidste ord.

\scene{\ensuremath{\emptyset} siger ingenting.}

\says{Clement} Anyways, lad os komme videre til næste spørgsmål.

\says{Clement} Hvad vil I gøre for tallenes repræsentation?

\scene{Spotlight på \ensuremath{\emptyset} og \ensuremath{\mathbb{N}}. Der går 10 sekunder, hvor ingen af dem svarer.}

\says{\ensuremath{\emptyset}} \ldots

\says{\ensuremath{\mathbb{N}}} \ldots

\says{Clement} Hvad fanden laver I?

\says{\ensuremath{\mathbb{N}}} Vi står og venter på, at den anden svarer.

\says{Clement} Nå, men nu er de 10 sekunder altså gået.

\says{Clement} \ensuremath{\mathbb{Z}} og \ensuremath{\mathbb{C}}, hvad vil I gøre for, at tallene bliver bedre repræsenteret?

\says{\ensuremath{\mathbb{Z}}} Vi mener jo, at vi som parti er store nok til at repræsentere hele samfundet.

\says{\ensuremath{\mathbb{Z}}} Vi er en partigruppe, som der sagtens kan tilføjes andre hele medlemmer til, uden at det hele bryder sammen.

\says{Clement} Men hvad så med de skæve individer?

\says{\ensuremath{\mathbb{Z}}} Dem holder vi uden for.

\says{\ensuremath{\mathbb{C}}} Det synes vi godt nok er lidt groft.

\says{\ensuremath{\mathbb{C}}} Hos os er der plads til både de imaginære og de skæve.

\says{\ensuremath{\mathbb{Z}}} Det må du have røget dig skæv for at tro.

\says{\ensuremath{\mathbb{Z}}} I har ikke engang en god orden i jeres parti.

\says{\ensuremath{\mathbb{C}}} Hov, hov du.

\says{\ensuremath{\mathbb{C}}} Det må du altså ikke sige, vi har altså følelser.

\says{Clement} Nu bryder jeg altså lige ind, fordi vi også gerne vil høre fra de to sidste partier.

\says{Clement} Der er mange, der siger, at I minder om hinanden.

\says{\ensuremath{\mathbb{Q}}} Det er jo rigtigt, at vi har haft et tæt forhold tidligere.

\says{Clement} Men?

\says{\ensuremath{\mathbb{Q}}} Well, orden er bare forsvundet fra deres parti.

\says{\ensuremath{\mathbb{R}}} Vi har da orden!

\says{\ensuremath{\mathbb{Q}}} Kun ved at tage mange mærkelige valg.

\scene{\ensuremath{\mathbb{N}_0} kommer ind med et tydeligt hjemmelavet podie og stiller sig mellem \ensuremath{\mathbb{N}} og \ensuremath{\mathbb{Z}}.}

\says{Clement} Vi går videre til næste spørgsmål.

\says{\ensuremath{\mathbb{N}_0}} Hey, hvad med mig?

\says{\ensuremath{\mathbb{N}}} Du er lille bitte parti, som ingen gider at høre på.

\says{\ensuremath{\mathbb{N}_0}} Det siger du bare, fordi jeg er dig, men bedre.

\says{Clement} Ninjaer, få ham ud.

\scene{Ninjaerne fjerner \ensuremath{\mathbb{N}_0} og hans podie}

\scene{Debatemne 3: Hvilke andre partier kunne I tænke jer at gå i forening med efter valget?}

\says{Clement} Hvilke andre partier kunne I tænke jer at gå i forening med efter valget?

\says{Clement} Vi starter hos jer, \ensuremath{\mathbb{R}}.

\says{\ensuremath{\mathbb{R}}} I alle de intervaller, vi har regeret. Det vil sige alle intervaller. Har vi kontinuerligt ført en politik, der fik os sikkert fra A til B uden nogen krumspring.

\says{\ensuremath{\mathbb{R}}} Modsat stort set alle de andre partier, som hele tiden hopper fra den ene holdning til den næste uden nogen sammenhæng.

\says{\ensuremath{\mathbb{R}}} Derfor så vi helst, at vi kunne danne et flertal med \ensuremath{\mathbb{C}}.

\says{\ensuremath{\mathbb{Z}}} Hvis man går sammen med jer, så kommer man jo bare til at føre jeres politik.

\says{\ensuremath{\mathbb{R}}} Netop!

\says{Clement} Nu er vi også blevet trætte af høre på jer to. Lad os vende til jer N.
Hvem vil egentlig gå sammen med, hvis jeres parti kommer ind i Folketinget?

\says{\ensuremath{\mathbb{N}}} Vi har gode positive holdninger.

\says{\ensuremath{\mathbb{N}}} Jeres holdninger er så ringe.

\says{\ensuremath{\mathbb{N}}} Så vi vil gerne arbejde alene.

\says{Clement} Hvad med jer, Liste \ensuremath{\emptyset}?

\says{\ensuremath{\mathbb{Z}}} I mener jo slet ikke noget.

\says{\ensuremath{\mathbb{Z}}} I vil altid i forening med alle.

\says{\ensuremath{\mathbb{Q}}} Vi må jo forholde os praktisk og rationelt til matematik.

\says{\ensuremath{\mathbb{Q}}} Derfor gider vi slet ikke arbejde sammen med sådan nogle små fringe-partier som \ensuremath{\mathbb{Z}} og \ensuremath{\mathbb{N}}.

\says{\ensuremath{\mathbb{N}}} Men vi er jo lige store?!

\scene{Tæppe}

\end{sketch}
\end{document}
